\answersfor{ch:sortbug}

\begin{solution}{ex:sortbug:1}
\emph{Correct value at 650\,nm.} The neighbours are $(600, 0.10)$ and $(700, 0.20)$.
\begin{align*}
t &= \frac{650 - 600}{700 - 600} = \frac{50}{100} = 0.5\\
y &= 0.10 + 0.5 \times (0.20 - 0.10)\\
  &= 0.10 + 0.05\\
  &= 0.15
\end{align*}
\emph{Listed 800, 700, 600.} The last listed sample is 600\,nm with value 0.10; the first is 800\,nm with value
0.30. Check 1: is $650 > 600$? Yes, so the last listed value is returned: $0.10$.

\emph{At 590\,nm.} Check 1: is $590 > 600$? No. Check 2: is $590 < 800$? Yes, so the first listed value is
returned: $0.30$. The correct answer is the held end value below the smallest sample, $0.10$. The descending list
gives the value at the \emph{other} end of the spectrum.
\end{solution}

\begin{solution}{ex:sortbug:2}
\emph{Correct row.} 620\,nm lies between 600 and 700\,nm:
\[
t = \frac{620 - 600}{700 - 600} = \frac{20}{100} = 0.2 .
\]
By $y = (1-t)\,y_{600} + t\,y_{700}$, the 700\,nm band gets weight $t = 0.2$ and the 600\,nm band $1 - t = 0.8$;
500\,nm gets 0. Row (listed order 700, 600, 500): $(0.2, 0.8, 0.0)$. Applied to $(0.40, 0.30, 0.10)$:
\begin{align*}
& 0.2 \times 0.40 + 0.8 \times 0.30 + 0 \times 0.10\\
&= 0.08 + 0.24 + 0\\
&= 0.32
\end{align*}
\emph{Unsorted row.} $620 > 500$, the last listed wavelength, so every column returns the last entry of its unit
vector: $(0, 0, 1)$. Applied to the pixel: $0 \times 0.40 + 0 \times 0.30 + 1 \times 0.10 = 0.10$.
\end{solution}

\begin{solution}{ex:sortbug:3}
\emph{Sort.} $\mathtt{wl} = (600, 800, 700)$. The smallest, 600, is at position 0; next, 700, at position 2; the
largest, 800, at position 1. So $\mathtt{argsort} = (0, 2, 1)$ and the sorted list is $(600, 700, 800)$.

\emph{Columns on the sorted list.} Grid 650 lies halfway between 600 and 700 ($t = 0.5$); grid 750 halfway between
700 and 800 ($t = 0.5$).
\begin{center}
\begin{tabular}{lccc}
\toprule
 & sorted 0 (600) & sorted 1 (700) & sorted 2 (800) \\
\midrule
650\,nm & 0.5 & 0.5 & 0 \\
750\,nm & 0 & 0.5 & 0.5 \\
\bottomrule
\end{tabular}
\end{center}
\emph{Write back.} Sorted column $i$ goes to original column $\mathtt{order}[i]$: sorted 0 $\to$ original 0
(600), sorted 1 $\to$ original 2 (700), sorted 2 $\to$ original 1 (800). In the listed order (600, 800, 700):
\[
\mat{M} = \begin{pmatrix} 0.5 & 0 & 0.5 \\ 0 & 0.5 & 0.5 \end{pmatrix}.
\]
\emph{Apply to $(1, 3, 2)$.}
\begin{align*}
\text{row 1:}\quad & 0.5 \times 1 + 0 \times 3 + 0.5 \times 2 = 0.5 + 0 + 1.0 = 1.5\\
\text{row 2:}\quad & 0 \times 1 + 0.5 \times 3 + 0.5 \times 2 = 0 + 1.5 + 1.0 = 2.5
\end{align*}
Check against the function itself: the pairs are $600 \to 1$, $700 \to 2$, $800 \to 3$, a straight line, so 650
gives 1.5 and 750 gives 2.5. The sort recovered the right answer from an unsorted list.
\end{solution}

\begin{solution}{ex:sortbug:4}
\emph{bench-v2.1.} With $p = 0.34$, $n = 2{,}817$:
\begin{align*}
p(1-p) &= 0.34 \times 0.66 = 0.2244\\
0.2244 / 2{,}817 &= 0.0000797\\
\sqrt{0.0000797} &= 0.0089 \quad\text{(0.89 points)}\\
35.1 - 33.1 &= 2.0 \text{ points}\\
2.0 / 0.89 &= 2.2 \text{ standard errors.}
\end{align*}
\emph{Real bench.} With $p = 0.30$, $n = 909$:
\begin{align*}
p(1-p) &= 0.30 \times 0.70 = 0.21\\
0.21 / 909 &= 0.000231\\
\sqrt{0.000231} &= 0.0152 \quad\text{(1.52 points)}\\
44.2 - 15.8 &= 28.4 \text{ points}\\
28.4 / 1.52 &= 18.7 \text{ standard errors.}
\end{align*}
The bug's effect is about eight times larger, in standard errors, than the largest model-to-model step on
bench-v2.1. A gap of nearly 19 standard errors between two readings of one model on the same scenes cannot be
noise; it can only be a difference in the data path.
\end{solution}

\begin{solution}{ex:sortbug:5}
The tests fed ascending wavelength lists, such as $[500, 600, 700]$. On an ascending list the precondition of
\texttt{np.interp} holds, so the sorted and unsorted code compute the same matrix and every test passes. The bug
only appears on real PRISMA files, which list their bands descending.

An assertion that makes the old code fail loudly, placed before the \texttt{np.interp} calls:
\begin{lstlisting}
assert np.all(np.diff(wl) > 0), "np.interp needs ascending wavelengths; sort wl first"
\end{lstlisting}
\texttt{np.diff(wl)} is the list of differences between neighbouring entries; all of them are positive exactly when
the list is strictly increasing. On $(700, 600, 500)$ the differences are $(-100, -100)$ and the assertion fails.
(The repository chose to sort instead of asserting, which also makes the function correct for any order;
\texttt{test\_descending\_band\_lists} pins that behaviour.)
\end{solution}
